\chapter{Unit 3: Fundamentals of Electrical Machines}

\section{Section A: Multiple Choice Question Practice (60 Curated Problems)}

\noindent\textbf{Q1. Mutual inductance between two coils is the property by which a changing current in one coil induces an EMF in the:} \hfill [Topic: mutual inductance and mutual coupling phenomena in transformer | Easy]
\begin{enumerate}[label=(\alph*),topsep=2pt,itemsep=1pt,partopsep=0pt,parsep=0pt]
    \item {Same coil only}
    \item {Iron core only}
    \item {Neighbouring coil}
    \item {Supply source}
\end{enumerate}
\vspace{2mm}

\noindent\textbf{Q2. The SI unit of mutual inductance is the:} \hfill [Topic: mutual inductance and mutual coupling phenomena in transformer | Easy]
\begin{enumerate}[label=(\alph*),topsep=2pt,itemsep=1pt,partopsep=0pt,parsep=0pt]
    \item {Farad}
    \item {Henry}
    \item {Weber}
    \item {Tesla}
\end{enumerate}
\vspace{2mm}

\noindent\textbf{Q3. In a transformer, magnetic coupling between the two windings is provided mainly by the:} \hfill [Topic: mutual inductance and mutual coupling phenomena in transformer | Easy]
\begin{enumerate}[label=(\alph*),topsep=2pt,itemsep=1pt,partopsep=0pt,parsep=0pt]
    \item {Laminated iron core}
    \item {Insulation layer}
    \item {Copper losses}
    \item {Air gap}
\end{enumerate}
\vspace{2mm}

\noindent\textbf{Q4. A transformer works on the principle of:} \hfill [Topic: transformer -- working | Easy]
\begin{enumerate}[label=(\alph*),topsep=2pt,itemsep=1pt,partopsep=0pt,parsep=0pt]
    \item {Mutual electromagnetic induction}
    \item {Self capacitance}
    \item {Chemical reaction}
    \item {Electrostatic induction}
\end{enumerate}
\vspace{2mm}

\noindent\textbf{Q5. A transformer can operate only with:} \hfill [Topic: transformer -- working | Easy]
\begin{enumerate}[label=(\alph*),topsep=2pt,itemsep=1pt,partopsep=0pt,parsep=0pt]
    \item {Direct current}
    \item {Static charge}
    \item {Alternating current}
    \item {Pulsating constant current}
\end{enumerate}
\vspace{2mm}

\noindent\textbf{Q6. A transformer that increases the output voltage is called a:} \hfill [Topic: transformer -- working | Easy]
\begin{enumerate}[label=(\alph*),topsep=2pt,itemsep=1pt,partopsep=0pt,parsep=0pt]
    \item {Current transformer}
    \item {Isolation transformer}
    \item {Step-down transformer}
    \item {Step-up transformer}
\end{enumerate}
\vspace{2mm}

\noindent\textbf{Q7. The turns ratio of a transformer is defined as the ratio of:} \hfill [Topic: concept of turns ratio and applications | Easy]
\begin{enumerate}[label=(\alph*),topsep=2pt,itemsep=1pt,partopsep=0pt,parsep=0pt]
    \item {Primary power to secondary power}
    \item {Primary turns to secondary turns}
    \item {Core area to window area}
    \item {Primary current to secondary current}
\end{enumerate}
\vspace{2mm}

\noindent\textbf{Q8. For an ideal transformer, the voltage ratio equals the turns ratio as V 2 V 1 = ?} \hfill [Topic: concept of turns ratio and applications | Easy]
\begin{enumerate}[label=(\alph*),topsep=2pt,itemsep=1pt,partopsep=0pt,parsep=0pt]
    \item {N 1 N 2 1}
    \item {N 2 N 1}
    \item {N 1 $\times$ N 2}
    \item {N 1 N 2}
\end{enumerate}
\vspace{2mm}

\noindent\textbf{Q9. In an ideal transformer, if the voltage is stepped up, the current in that winding is:} \hfill [Topic: concept of turns ratio and applications | Easy]
\begin{enumerate}[label=(\alph*),topsep=2pt,itemsep=1pt,partopsep=0pt,parsep=0pt]
    \item {Reduced to zero}
    \item {Unchanged}
    \item {Stepped up equally}
    \item {Stepped down}
\end{enumerate}
\vspace{2mm}

\noindent\textbf{Q10. The device used to convert the alternating EMF in a DC generator armature into DC at the output is the:} \hfill [Topic: dc machines- working principles, classification | Easy]
\begin{enumerate}[label=(\alph*),topsep=2pt,itemsep=1pt,partopsep=0pt,parsep=0pt]
    \item {Transformer}
    \item {Commutator}
    \item {Slip ring}
    \item {Capacitor}
\end{enumerate}
\vspace{2mm}

\noindent\textbf{Q11. Based on the field connection, DC machines are commonly classified as series, shunt and:} \hfill [Topic: dc machines- working principles, classification | Easy]
\begin{enumerate}[label=(\alph*),topsep=2pt,itemsep=1pt,partopsep=0pt,parsep=0pt]
    \item {Compound}
    \item {Synchronous}
    \item {Parallel}
    \item {Induction}
\end{enumerate}
\vspace{2mm}

\noindent\textbf{Q12. A DC motor converts electrical energy into:} \hfill [Topic: dc machines- working principles, classification | Easy]
\begin{enumerate}[label=(\alph*),topsep=2pt,itemsep=1pt,partopsep=0pt,parsep=0pt]
    \item {Light energy}
    \item {Chemical energy}
    \item {Mechanical energy}
    \item {Heat energy only}
\end{enumerate}
\vspace{2mm}

\noindent\textbf{Q13. A starter is used with a DC motor mainly to limit the:} \hfill [Topic: starting, speed control and applications of dc motors | Easy]
\begin{enumerate}[label=(\alph*),topsep=2pt,itemsep=1pt,partopsep=0pt,parsep=0pt]
    \item {Field flux}
    \item {Running speed}
    \item {Output torque}
    \item {High starting current}
\end{enumerate}
\vspace{2mm}

\noindent\textbf{Q14. Which DC motor is best suited for applications requiring very high starting torque, such as traction?} \hfill [Topic: starting, speed control and applications of dc motors | Easy]
\begin{enumerate}[label=(\alph*),topsep=2pt,itemsep=1pt,partopsep=0pt,parsep=0pt]
    \item {Series motor}
    \item {Shunt motor}
    \item {Synchronous motor}
    \item {Separately excited motor}
\end{enumerate}
\vspace{2mm}

\noindent\textbf{Q15. The rotating magnetic field in a three-phase induction motor is produced by the:} \hfill [Topic: working principle of single and three phase induction motors | Easy]
\begin{enumerate}[label=(\alph*),topsep=2pt,itemsep=1pt,partopsep=0pt,parsep=0pt]
    \item {Three-phase stator windings}
    \item {Permanent magnets}
    \item {Commutator}
    \item {DC field supply}
\end{enumerate}
\vspace{2mm}

\noindent\textbf{Q16. A single-phase induction motor is not self-starting because its stator produces a:} \hfill [Topic: working principle of single and three phase induction motors | Easy]
\begin{enumerate}[label=(\alph*),topsep=2pt,itemsep=1pt,partopsep=0pt,parsep=0pt]
    \item {Rotating magnetic field}
    \item {Zero magnetic field}
    \item {Constant DC field}
    \item {Pulsating magnetic field}
\end{enumerate}
\vspace{2mm}

\noindent\textbf{Q17. Three-phase induction motors are widely used in industry mainly because they are:} \hfill [Topic: applications of ac motors | Easy]
\begin{enumerate}[label=(\alph*),topsep=2pt,itemsep=1pt,partopsep=0pt,parsep=0pt]
    \item {Expensive and complex}
    \item {Rugged and low maintenance}
    \item {Suitable for DC supply}
    \item {Unable to run on load}
\end{enumerate}
\vspace{2mm}

\noindent\textbf{Q18. Fleming's left-hand rule is used to find the direction of force in a:} \hfill [Topic: Fleming's left hand and right hand rule | Easy]
\begin{enumerate}[label=(\alph*),topsep=2pt,itemsep=1pt,partopsep=0pt,parsep=0pt]
    \item {Transformer}
    \item {Generator}
    \item {Rectifier}
    \item {Motor}
\end{enumerate}
\vspace{2mm}

\noindent\textbf{Q19. In Fleming's right-hand rule, the thumb represents the direction of:} \hfill [Topic: Fleming's left hand and right hand rule | Easy]
\begin{enumerate}[label=(\alph*),topsep=2pt,itemsep=1pt,partopsep=0pt,parsep=0pt]
    \item {Motion of the conductor}
    \item {Magnetic field}
    \item {Applied force}
    \item {Induced current}
\end{enumerate}
\vspace{2mm}

\noindent\textbf{Q20. If a DC voltage is applied to a transformer primary, the output voltage at the secondary will be:} \hfill [Topic: transformer on DC | Easy]
\begin{enumerate}[label=(\alph*),topsep=2pt,itemsep=1pt,partopsep=0pt,parsep=0pt]
    \item {Zero in steady state}
    \item {Negative of input}
    \item {Doubled}
    \item {Equal to input}
\end{enumerate}
\vspace{2mm}

\noindent\textbf{Q21. A single-phase transformer has 400 primary turns and 100 secondary turns. If the primary voltage is 230 V , what is the secondary voltage (assuming ideal transformer)?} \hfill [Topic: concept of turns ratio and applications | Medium]
\begin{enumerate}[label=(\alph*),topsep=2pt,itemsep=1pt,partopsep=0pt,parsep=0pt]
    \item {460 V}
    \item {57.5 V}
    \item {115 V}
    \item {920 V}
\end{enumerate}
\vspace{2mm}

\noindent\textbf{Q22. An ideal step-up transformer increases voltage by a factor of 5. If the primary current is 10 A , the secondary current is:} \hfill [Topic: concept of turns ratio and applications | Medium]
\begin{enumerate}[label=(\alph*),topsep=2pt,itemsep=1pt,partopsep=0pt,parsep=0pt]
    \item {50 A}
    \item {10 A}
    \item {2 A}
    \item {5 A}
\end{enumerate}
\vspace{2mm}

\noindent\textbf{Q23. What happens when a DC supply is connected to the primary of a transformer?} \hfill [Topic: transformer on DC | Medium]
\begin{enumerate}[label=(\alph*),topsep=2pt,itemsep=1pt,partopsep=0pt,parsep=0pt]
    \item {The secondary current reverses direction periodically}
    \item {The transformer works normally with slightly reduced efficiency}
    \item {A steady voltage appears across the secondary equal to the turns ratio times the input}
    \item {The constant flux induces no EMF in the secondary, and the low winding resistance causes a very large steady current that can overheat and burn the primary coil}
\end{enumerate}
\vspace{2mm}

\noindent\textbf{Q24. Two coils have a mutual inductance of 0.2 H . If the current in one coil changes at a rate of 500 A/s , the EMF induced in the other coil is:} \hfill [Topic: mutual inductance and mutual coupling phenomena in transformer | Medium]
\begin{enumerate}[label=(\alph*),topsep=2pt,itemsep=1pt,partopsep=0pt,parsep=0pt]
    \item {10 V}
    \item {2500 V}
    \item {250 V}
    \item {100 V}
\end{enumerate}
\vspace{2mm}

\noindent\textbf{Q25. Two coils with self-inductances L 1 = 0.1 H and L 2 = 0.4 H are perfectly coupled ( k = 1 ). The mutual inductance is:} \hfill [Topic: mutual inductance and mutual coupling phenomena in transformer | Medium]
\begin{enumerate}[label=(\alph*),topsep=2pt,itemsep=1pt,partopsep=0pt,parsep=0pt]
    \item {0.2 H}
    \item {0.1 H}
    \item {0.25 H}
    \item {0.5 H}
\end{enumerate}
\vspace{2mm}

\noindent\textbf{Q26. The primary function of the laminated iron core in a transformer is to:} \hfill [Topic: transformer -- working | Medium]
\begin{enumerate}[label=(\alph*),topsep=2pt,itemsep=1pt,partopsep=0pt,parsep=0pt]
    \item {Convert AC into pulsating DC}
    \item {Provide a low-reluctance path for the magnetic flux while reducing eddy current losses}
    \item {Store electrical charge between the windings}
    \item {Increase the resistance of the windings}
\end{enumerate}
\vspace{2mm}

\noindent\textbf{Q27. The EMF equation of a transformer is E = 4.44 f N $\phi$ m . If the supply frequency is doubled while keeping voltage and turns constant, the maximum flux $\phi$ m will:} \hfill [Topic: transformer -- working | Medium]
\begin{enumerate}[label=(\alph*),topsep=2pt,itemsep=1pt,partopsep=0pt,parsep=0pt]
    \item {Reduce to half}
    \item {Double}
    \item {Reduce to one-fourth}
    \item {Remain the same}
\end{enumerate}
\vspace{2mm}

\noindent\textbf{Q28. Compared to a two-winding transformer of the same rating, an auto-transformer requires:} \hfill [Topic: auto-transformer | Medium]
\begin{enumerate}[label=(\alph*),topsep=2pt,itemsep=1pt,partopsep=0pt,parsep=0pt]
    \item {The same amount of copper but larger core}
    \item {Less copper and has higher efficiency, but offers no electrical isolation between primary and secondary}
    \item {More copper and has lower efficiency}
    \item {Two completely separate windings with full isolation}
\end{enumerate}
\vspace{2mm}

\noindent\textbf{Q29. A current transformer (CT) is used to measure large currents. During operation, its secondary should:} \hfill [Topic: instrument transformers | Medium]
\begin{enumerate}[label=(\alph*),topsep=2pt,itemsep=1pt,partopsep=0pt,parsep=0pt]
    \item {Never be left open-circuited, as dangerously high voltages can be induced}
    \item {Be connected to the supply voltage directly}
    \item {Be connected across a high-value resistor only}
    \item {Always be left open-circuited for safety}
\end{enumerate}
\vspace{2mm}

\noindent\textbf{Q30. A potential transformer (PT) with ratio 11000/110 V shows 100 V on its voltmeter. The actual line voltage is:} \hfill [Topic: instrument transformers | Medium]
\begin{enumerate}[label=(\alph*),topsep=2pt,itemsep=1pt,partopsep=0pt,parsep=0pt]
    \item {1100 V}
    \item {1000 V}
    \item {10000 V}
    \item {11000 V}
\end{enumerate}
\vspace{2mm}

\noindent\textbf{Q31. Fleming's left-hand rule is used to determine the direction of:} \hfill [Topic: Fleming's left hand and right hand rule | Medium]
\begin{enumerate}[label=(\alph*),topsep=2pt,itemsep=1pt,partopsep=0pt,parsep=0pt]
    \item {Force on a current-carrying conductor in a magnetic field (motor action)}
    \item {Induced EMF in a generator}
    \item {Magnetic flux inside a solenoid}
    \item {The resistance of a conductor}
\end{enumerate}
\vspace{2mm}

\noindent\textbf{Q32. In Fleming's right-hand rule, the forefinger, thumb, and middle finger respectively represent:} \hfill [Topic: Fleming's left hand and right hand rule | Medium]
\begin{enumerate}[label=(\alph*),topsep=2pt,itemsep=1pt,partopsep=0pt,parsep=0pt]
    \item {Current, field, and motion}
    \item {Motion, current, and field}
    \item {Field, motion, and induced current}
    \item {Field, current, and force}
\end{enumerate}
\vspace{2mm}

\noindent\textbf{Q33. In a DC motor, the back EMF ( E b ) at the instant of starting is:} \hfill [Topic: dc machines- working principles, classification, starting, speed control and applications of dc motors | Medium]
\begin{enumerate}[label=(\alph*),topsep=2pt,itemsep=1pt,partopsep=0pt,parsep=0pt]
    \item {Maximum, equal to the supply voltage}
    \item {Zero, because the armature is not yet rotating}
    \item {Negative and opposing the current}
    \item {Equal to half the supply voltage}
\end{enumerate}
\vspace{2mm}

\noindent\textbf{Q34. A DC shunt motor runs at 1000 rpm . If the flux is reduced to 80\% of its original value (armature EMF roughly constant), the speed becomes approximately:} \hfill [Topic: dc machines- working principles, classification, starting, speed control and applications of dc motors | Medium]
\begin{enumerate}[label=(\alph*),topsep=2pt,itemsep=1pt,partopsep=0pt,parsep=0pt]
    \item {1250 rpm}
    \item {640 rpm}
    \item {800 rpm}
    \item {1000 rpm}
\end{enumerate}
\vspace{2mm}

\noindent\textbf{Q35. Which DC motor is best suited for applications requiring high starting torque, such as electric traction and cranes?} \hfill [Topic: dc machines- working principles, classification, starting, speed control and applications of dc motors | Medium]
\begin{enumerate}[label=(\alph*),topsep=2pt,itemsep=1pt,partopsep=0pt,parsep=0pt]
    \item {DC series motor}
    \item {DC compound (cumulative) motor}
    \item {Separately excited DC motor}
    \item {DC shunt motor}
\end{enumerate}
\vspace{2mm}

\noindent\textbf{Q36. A DC series motor should never be started without a mechanical load because:} \hfill [Topic: dc machines- working principles, classification, starting, speed control and applications of dc motors | Medium]
\begin{enumerate}[label=(\alph*),topsep=2pt,itemsep=1pt,partopsep=0pt,parsep=0pt]
    \item {The armature current becomes zero}
    \item {At no load the flux is very low, causing the speed to rise dangerously high and risk mechanical damage}
    \item {The commutator will stop working}
    \item {It will fail to develop any torque at all}
\end{enumerate}
\vspace{2mm}

\noindent\textbf{Q37. A 4-pole, three-phase induction motor is connected to a 50 Hz supply. Its synchronous speed is:} \hfill [Topic: working principle of single and three phase induction motors | Medium]
\begin{enumerate}[label=(\alph*),topsep=2pt,itemsep=1pt,partopsep=0pt,parsep=0pt]
    \item {1000 rpm}
    \item {750 rpm}
    \item {1500 rpm}
    \item {3000 rpm}
\end{enumerate}
\vspace{2mm}

\noindent\textbf{Q38. A three-phase induction motor has a synchronous speed of 1500 rpm and runs at 1440 rpm . The percentage slip is:} \hfill [Topic: working principle of single and three phase induction motors | Medium]
\begin{enumerate}[label=(\alph*),topsep=2pt,itemsep=1pt,partopsep=0pt,parsep=0pt]
    \item {10\%}
    \item {4\%}
    \item {6\%}
    \item {2\%}
\end{enumerate}
\vspace{2mm}

\noindent\textbf{Q39. A single-phase induction motor is not self-starting because:} \hfill [Topic: working principle of single and three phase induction motors | Medium]
\begin{enumerate}[label=(\alph*),topsep=2pt,itemsep=1pt,partopsep=0pt,parsep=0pt]
    \item {The supply frequency is too low to start it}
    \item {It has no rotor conductors}
    \item {The rotor is made of a permanent magnet}
    \item {The single-phase winding produces a pulsating field, giving zero net starting torque}
\end{enumerate}
\vspace{2mm}

\noindent\textbf{Q40. Which AC motor is most commonly used in fans, pumps, and household appliances due to its ruggedness, low cost, and low maintenance?} \hfill [Topic: applications of ac motors | Medium]
\begin{enumerate}[label=(\alph*),topsep=2pt,itemsep=1pt,partopsep=0pt,parsep=0pt]
    \item {Synchronous motor}
    \item {Wound-rotor induction motor}
    \item {Universal motor}
    \item {Squirrel-cage induction motor}
\end{enumerate}
\vspace{2mm}

\noindent\textbf{Q41. Two coils have self-inductances L 1 = 0.4 H and L 2 = 0.9 H . When connected in series aiding, the total inductance is 1.78 H . What is the coefficient of coupling k ?} \hfill [Topic: mutual inductance and mutual coupling phenomena in transformer | Hard]
\begin{enumerate}[label=(\alph*),topsep=2pt,itemsep=1pt,partopsep=0pt,parsep=0pt]
    \item {0.80}
    \item {0.65}
    \item {0.40}
    \item {0.50}
\end{enumerate}
\vspace{2mm}

\noindent\textbf{Q42. A transformer rated for 230 V AC is accidentally connected to a 230 V DC supply. Which outcome best describes what happens and why?} \hfill [Topic: transformer on DC | Hard]
\begin{enumerate}[label=(\alph*),topsep=2pt,itemsep=1pt,partopsep=0pt,parsep=0pt]
    \item {The primary draws a very large current limited only by its small winding resistance, risking burnout, because DC produces no back-EMF of self-induction in steady state and the inductive reactance is zero}
    \item {The core saturates but current stays at rated value}
    \item {The secondary delivers 230 V DC safely}
    \item {The transformer steps the DC voltage up or down normally}
\end{enumerate}
\vspace{2mm}

\noindent\textbf{Q43. A step-down transformer has turns ratio 10 : 1 . The secondary is loaded with 5 $\Omega$ . What impedance is reflected to the primary?} \hfill [Topic: concept of turns ratio and applications | Hard]
\begin{enumerate}[label=(\alph*),topsep=2pt,itemsep=1pt,partopsep=0pt,parsep=0pt]
    \item {250 $\Omega$}
    \item {0.05 $\Omega$}
    \item {500 $\Omega$}
    \item {50 $\Omega$}
\end{enumerate}
\vspace{2mm}

\noindent\textbf{Q44. In a practical transformer at no-load, the no-load current lags the applied voltage by nearly 8 0 $^\circ$ . What does this large angle primarily indicate?} \hfill [Topic: transformer -- working | Hard]
\begin{enumerate}[label=(\alph*),topsep=2pt,itemsep=1pt,partopsep=0pt,parsep=0pt]
    \item {The core losses are extremely high}
    \item {The magnetizing (reactive) component dominates over the small core-loss (active) component}
    \item {The transformer is heavily loaded}
    \item {The winding resistance is dominant}
\end{enumerate}
\vspace{2mm}

\noindent\textbf{Q45. A 1 kV A , 200/100 V step-down transformer is reconnected as an auto-transformer to give 200/300 V . What is the new kVA rating as an auto-transformer?} \hfill [Topic: auto-transformer | Hard]
\begin{enumerate}[label=(\alph*),topsep=2pt,itemsep=1pt,partopsep=0pt,parsep=0pt]
    \item {1 kV A}
    \item {2 kV A}
    \item {1.5 kV A}
    \item {3 kV A}
\end{enumerate}
\vspace{2mm}

\noindent\textbf{Q46. Compared to an equivalent two-winding transformer of the same output, an auto-transformer transfers part of the power conductively. For a turns ratio where N 1 N 2 = 0.8 (step-down), what fraction of the total power is transferred inductively (transformed)?} \hfill [Topic: auto-transformer | Hard]
\begin{enumerate}[label=(\alph*),topsep=2pt,itemsep=1pt,partopsep=0pt,parsep=0pt]
    \item {1.0}
    \item {0.2}
    \item {0.5}
    \item {0.8}
\end{enumerate}
\vspace{2mm}

\noindent\textbf{Q47. Why should the secondary of a current transformer (CT) never be left open-circuited while the primary carries current?} \hfill [Topic: instrument transformers | Hard]
\begin{enumerate}[label=(\alph*),topsep=2pt,itemsep=1pt,partopsep=0pt,parsep=0pt]
    \item {The absence of secondary demagnetizing MMT lets the full primary MMT magnetize the core, producing dangerously high flux, core overheating and a very high secondary voltage}
    \item {The secondary voltage drops to zero}
    \item {The CT ratio simply becomes inaccurate but stays safe}
    \item {The primary current falls to zero immediately}
\end{enumerate}
\vspace{2mm}

\noindent\textbf{Q48. A potential transformer (PT) is used to measure a 11 kV line voltage with a standard 110 V voltmeter. If the measured ratio error causes the voltmeter to read 108 V for an actual 11 kV , what is the approximate percentage ratio error?} \hfill [Topic: instrument transformers | Hard]
\begin{enumerate}[label=(\alph*),topsep=2pt,itemsep=1pt,partopsep=0pt,parsep=0pt]
    \item {+ 2.0\%}
    \item {+ 1.8\%}
    \item {- 2.0\%}
    \item {- 1.8\%}
\end{enumerate}
\vspace{2mm}

\noindent\textbf{Q49. In a DC generator, the function of the commutator is best described as:} \hfill [Topic: dc machines- working principles | Hard]
\begin{enumerate}[label=(\alph*),topsep=2pt,itemsep=1pt,partopsep=0pt,parsep=0pt]
    \item {A mechanical rectifier that converts the internally generated alternating EMF into unidirectional (DC) output at the brushes}
    \item {A device that increases the generated EMF}
    \item {A field-weakening control element}
    \item {A frequency converter for the AC output}
\end{enumerate}
\vspace{2mm}

\noindent\textbf{Q50. A DC shunt motor runs at 1000 r p m drawing an armature current of 20 A from a 220 V supply with armature resistance 0.5 $\Omega$ . If flux is held constant and the load torque is unchanged, what series armature resistance must be added to reduce the speed to 800 r p m ?} \hfill [Topic: speed control and applications of dc motors | Hard]
\begin{enumerate}[label=(\alph*),topsep=2pt,itemsep=1pt,partopsep=0pt,parsep=0pt]
    \item {2.1 $\Omega$}
    \item {3.0 $\Omega$}
    \item {1.5 $\Omega$}
    \item {0.5 $\Omega$}
\end{enumerate}
\vspace{2mm}

\noindent\textbf{Q51. A 250 V DC shunt motor has armature resistance 0.4 $\Omega$ . To limit the starting current to twice the full-load current of 40 A , what starter resistance is required at the instant of starting?} \hfill [Topic: starting | Hard]
\begin{enumerate}[label=(\alph*),topsep=2pt,itemsep=1pt,partopsep=0pt,parsep=0pt]
    \item {2.5 $\Omega$}
    \item {2.725 $\Omega$}
    \item {6.25 $\Omega$}
    \item {3.125 $\Omega$}
\end{enumerate}
\vspace{2mm}

\noindent\textbf{Q52. In a DC series motor, if the load torque suddenly drops to nearly zero (no-load), what happens and why?} \hfill [Topic: classification | Hard]
\begin{enumerate}[label=(\alph*),topsep=2pt,itemsep=1pt,partopsep=0pt,parsep=0pt]
    \item {The speed drops to zero}
    \item {The speed stays constant due to flux regulation}
    \item {The motor stalls immediately}
    \item {The speed rises dangerously high (racing) because the flux falls with the small current, and since N $\propto$ 1/ $\phi$ at fixed voltage the speed increases without bound}
\end{enumerate}
\vspace{2mm}

\noindent\textbf{Q53. A 4 -pole, 50 Hz three-phase induction motor runs at 1440 r p m . What is the frequency of the rotor-induced EMF?} \hfill [Topic: working principle of single and three phase induction motors | Hard]
\begin{enumerate}[label=(\alph*),topsep=2pt,itemsep=1pt,partopsep=0pt,parsep=0pt]
    \item {4 Hz}
    \item {2 Hz}
    \item {48 Hz}
    \item {50 Hz}
\end{enumerate}
\vspace{2mm}

\noindent\textbf{Q54. Why is a single-phase induction motor not self-starting, according to the double-field revolving theory?} \hfill [Topic: working principle of single and three phase induction motors | Hard]
\begin{enumerate}[label=(\alph*),topsep=2pt,itemsep=1pt,partopsep=0pt,parsep=0pt]
    \item {The pulsating field splits into two equal and opposite rotating fields that produce equal and opposite torques, giving zero net starting torque at standstill}
    \item {The rotor bars are open-circuited at start}
    \item {The stator produces no magnetic field at start}
    \item {The supply frequency is too low to induce rotor current}
\end{enumerate}
\vspace{2mm}

\noindent\textbf{Q55. A conductor carries current in the + x direction and lies in a magnetic field pointing in the + y direction. Using Fleming's left-hand rule, in which direction is the force on the conductor?} \hfill [Topic: Fleming's left hand and right hand rule | Hard]
\begin{enumerate}[label=(\alph*),topsep=2pt,itemsep=1pt,partopsep=0pt,parsep=0pt]
    \item {+ z direction}
    \item {- x direction}
    \item {- z direction}
    \item {+ y direction}
\end{enumerate}
\vspace{2mm}

\noindent\textbf{Q56. Which statement correctly distinguishes the two Fleming's rules in the context of electrical machines?} \hfill [Topic: Fleming's left hand and right hand rule | Hard]
\begin{enumerate}[label=(\alph*),topsep=2pt,itemsep=1pt,partopsep=0pt,parsep=0pt]
    \item {Both rules give force direction, differing only by hand}
    \item {Both rules apply only to transformers, not rotating machines}
    \item {The right-hand rule gives induced EMF/current direction in a generator; the left-hand rule gives force direction in a motor}
    \item {The left-hand rule applies to generators; the right-hand rule applies to motors}
\end{enumerate}
\vspace{2mm}

\noindent\textbf{Q57. For a constant-speed drive such as a large industrial fan or pump where power factor correction is also desired, which AC motor is most suitable?} \hfill [Topic: applications of ac motors | Hard]
\begin{enumerate}[label=(\alph*),topsep=2pt,itemsep=1pt,partopsep=0pt,parsep=0pt]
    \item {Capacitor-start induction motor}
    \item {Universal motor}
    \item {Synchronous motor, since it runs at constant speed and can be over-excited to deliver leading reactive power}
    \item {Single-phase shaded-pole motor}
\end{enumerate}
\vspace{2mm}

\noindent\textbf{Q58. A transformer has iron loss 400 W and full-load copper loss 400 W . At what fraction of full load does maximum efficiency occur, and how does the copper loss there compare with the iron loss?} \hfill [Topic: transformer -- working | Hard]
\begin{enumerate}[label=(\alph*),topsep=2pt,itemsep=1pt,partopsep=0pt,parsep=0pt]
    \item {At full load; copper loss is double the iron loss}
    \item {At half load; copper loss equals 200 W}
    \item {At quarter load; copper loss equals iron loss}
    \item {At full load; copper loss equals iron loss ( 400 W each)}
\end{enumerate}
\vspace{2mm}

\noindent\textbf{Q59. In field (flux) control of a DC shunt motor above base speed, the output power is approximately constant. If flux is reduced to 0.5 times its original value with armature voltage and current held at rated, the torque capability becomes:} \hfill [Topic: speed control and applications of dc motors | Hard]
\begin{enumerate}[label=(\alph*),topsep=2pt,itemsep=1pt,partopsep=0pt,parsep=0pt]
    \item {Halved, since torque $\propto$ $\phi$ I a while speed roughly doubles keeping power constant}
    \item {Reduced to one quarter}
    \item {Unchanged}
    \item {Doubled}
\end{enumerate}
\vspace{2mm}

\noindent\textbf{Q60. Two magnetically coupled coils connected in series carry the same current. When connected series-aiding the total inductance is 2.2 H and when series-opposing it is 1.0 H . What is the mutual inductance M ?} \hfill [Topic: mutual inductance and mutual coupling phenomena in transformer | Hard]
\begin{enumerate}[label=(\alph*),topsep=2pt,itemsep=1pt,partopsep=0pt,parsep=0pt]
    \item {0.15 H}
    \item {1.2 H}
    \item {0.3 H}
    \item {0.6 H}
\end{enumerate}
\vspace{2mm}


\newpage
\section{Section B: Comprehensive Analytical Problems with Step-by-Step Solutions}

\begin{subjectiveq}{3}{1}
A $25\text{ kVA}$, $2200/220\text{ V}$, $50\text{ Hz}$ single-phase distribution transformer has iron loss of $350\text{ W}$ and full-load copper loss of $400\text{ W}$. Calculate:\begin{enumerate}[label=(\alph*)]\item The efficiency at full-load and $0.8$ power factor lagging.\item The load fraction $x$ at which maximum efficiency occurs, and the value of maximum efficiency at unity power factor.\item Explain why transformers cannot operate on a steady D.C. supply.\end{enumerate}
\end{subjectiveq}

\begin{solbox}
\textbf{Step 1: Efficiency at Full Load ($x = 1.0, \cos\phi = 0.8$)}\\Output power: $P_{out} = S \times \cos\phi = 25000 \times 0.8 = 20000\text{ W} = 20\text{ kW}$.Total losses: $P_{losses} = P_i + P_{cu} = 350\text{ W} + 400\text{ W} = 750\text{ W} = 0.75\text{ kW}$.\[ \eta_{FL} = \frac{20000}{20000 + 750} \times 100\% \approx 96.39\% \]\textbf{Step 2: Maximum Efficiency Condition}\\Maximum efficiency occurs when copper loss equals iron loss ($x^2 P_{cu} = P_i$):\[ x = \sqrt{\frac{P_i}{P_{cu}}} = \sqrt{\frac{350}{400}} \approx 0.9354 \text{ (at } 93.54\% \text{ of rated load)} \]At this fraction and unity power factor ($\cos\phi = 1.0$):\[ P_{out} = 0.9354 \times 25000 \times 1.0 = 23385\text{ W} \]Total losses: $P_{losses} = 2 P_i = 700\text{ W}$.\[ \eta_{max} = \frac{23385}{23385 + 700} \times 100\% \approx 97.09\% \]\textbf{Step 3: Phenomenon of Transformer on Steady D.C.}\\Faraday's Law states induced EMF is $e = -N \frac{d\Phi}{dt}$. With steady D.C., $\frac{d\Phi}{dt} = 0$, so no counter-EMF is generated in the primary winding. The primary current is limited solely by low winding resistance $R_1$: $I_{dc} = V_{dc}/R_1$. This draws massive inrush current, causing immediate core saturation, extreme Joule heating, insulation breakdown, and transformer destruction.
\end{solbox}

\begin{examinertip}
\textbf{Marking Allocation}: Full load calculation: \textbf{3.5 Marks}; Condition for maximum efficiency: \textbf{3.5 Marks}; DC saturation explanation: \textbf{3.0 Marks}.
\end{examinertip}
\vspace{4mm}
